\section{The selected detector scales and Laurent circuit}

This section concerns the selected fifth-round construction. Its builder is
\texttt{glazer\_fifth\_reduction\_package/construction/builder.cjs}, SHA256
\texttt{68cb3cb1a809766aa102b8c43444d8b38c02760751e773a960489c327e33c12b}.
The fixed coefficient is $S_0=3^{2^{13}}$, not an extra quantified variable.
The earlier source is pinned at commit
\texttt{036cebe4f1a4aa7ca89820a9b48f3629b4c364d6}.
The uniform arithmetic inputs remain those stated in the accompanying
arithmetic proof; the finite source checks below do not independently
establish those inputs.

\subsection{One fewer coordinate for the third height scale}

Write the four-root abscissa product as
\[
 C=\prod_{i=1}^4x(Q_i)=C_0+C_1\rho,\qquad\rho^2=2.
\]
The selected rational scales are
\[
 1,\qquad C_0^2,\qquad(C_0+C_1)^2.
\]
The third replaces $(1+C_1^2)^2$ and removes its inner-square coordinate.
The physical norm anchor still uses $C_0^2$ and is unchanged by this
particular simplification.

At both real embeddings the root equations $x=h^2$ and
$j^2=x^2-(3-\rho)^2$ give $x\ge3-\sigma(\rho)>1$. Consequently
\[
 C_+\ge193-132\rho,\qquad C_-\ge193+132\rho.
\]
The identity
\[
 C_0+C_1=\frac{1+1/\rho}{2}C_+
         +\frac{1-1/\rho}{2}C_-
\]
has positive weights. It gives $C_0+C_1\ge61$, while $C_0\ge193$.
Thus every selected scale is a positive nonzero rational square and is
at least one. These are algebraic-domain facts on every source solution;
they do not depend on a rank or height statement.

The actual product-height proof is in
\texttt{docs/quartic-523-540/addition.tex}, lines296--350. For the new
fractional-linear coordinate $\phi(t)=-49+1/t$, the elementary height
comparison is $h(t)\le h(\phi(t))+\log50$. Therefore if a nonzero integer
$N$ is divisible by the $k$th powers of the three contact radicals, put
\[
 T=(H+\log50)|N|^{19/k}.
\]
The detector height estimates give
$h(\lambda),h(C_0^2\lambda),h((C_0+C_1)^2\lambda)\le T$.
Height subadditivity and $h(z^2)=2h(z)$ give
\[
 h(C_0),h(C_0+C_1)\le T,\qquad
 h(C_1)\le2T+\log2.
\]
Hence
\[
 h(C)\le h(C_0)+h(C_1)+\tfrac32\log2
       \le3T+\tfrac52\log2.
\]
The selected height constant satisfies $H\ge4$. Since
\[
 3\log50+\tfrac52\log2<16\le4H
\]
and $|N|^{19/k}\ge1$, the same retained bound follows:
\[
 h(C)<7H|N|^{19/k},\qquad B^2-6<7H|N|^{19/k}.
\]
For the numerical inequality it is already enough to use
$\log50<4$ and $\log2<7/10$. The intermediate height estimate is weaker
than its predecessor, but the retained constant $7H$ has sufficient slack.
This checks the changed coordinate and the changed scale together; the
old $\log3$ error is not used for the new coordinate.

In completeness the roots and scales are fixed before prime selection.
The source chooses $D_0$ such that every $D_0^2\theta/A$ is integral and
excludes the prime factors of the fixed scales and contact differences.
Its required scale hypothesis is a positive nonzero rational square,
which the replacement satisfies. The actual passage
\texttt{manuscripts/quartic\_v10.tex}, lines11244--11257, explicitly
deletes repeated contact forms while retaining their incidences in all
detectors. Thus a possible equality of the second and third scales is
harmless. The five forms within any one detector remain distinct because
its scale is nonzero. No new distinctness assumption is introduced.

\subsection{Selected contacts, positivity and the guard}

The coordinate and guard are
\[
 \phi(t)=-49+1/t,\qquad \delta=98t-1,
\]
with
\[
 \Theta=\{0,1/49,1/50,1/2450,32/1225\},\qquad
 F(t)=\prod_{\theta\in\Theta}(t-\theta).
\]
The five images under $\phi$ are respectively
$\infty,0,1,2401,-343/32$. The matrix of $\phi$ has determinant $-1$.
Contact denominators and pairwise differences have support in
$\{2,3,5,7\}$. The values of $\delta$ on the contacts are
\[
 -1,\quad1,\quad24/25,\quad-24/25,\quad39/25,
\]
which introduce only the additional prime13. The precise remaining
arithmetic consumers are checked in \texttt{arithmetic\_round5.tex};
the selected exceptional set is $\{2,3,5,7,13,17\}$.

The source imposes
\[
 \lambda=32/1225+u^2+v^2.
\]
Since $A\ge1$, $t=A\lambda$ is at least the largest contact. Therefore
$\delta\ge39/25>0$ and $F(t)\ge0$. The equation
\[
 eF(t)=(A\delta)^5
\]
forces $F(t)>0$ and $e>0$. In particular all Laurent denominators used
below are nonzero before invoking any pole equation.

The extra positivity equation has canonical completeness witnesses.
The unit-scale contact $32/1225$ is retained among the prime forms, so
the inherited prime selection gives
$\lambda-32/1225=p/(D_0w)^2$ with $p$ the norm of a selected Gaussian
prime. Dividing its two integer coefficients by $D_0w$ gives the two
rational square witnesses. Repeated-form merging retains this incidence.
Thus the extra equation uses the existing Gaussian norm conclusion,
not a new assertion that every positive rational is a sum of two squares.

\subsection{A same-cost order-forty-four chain}

For each scale set
\[
 e=\frac{(A\delta)^5}{F(t)},\qquad
 q_1=\frac e{\delta^3},\qquad
 R=(e\delta)^{-1}.
\]
Use the six intermediate powers
\[
 e^{-1},\ e^{-3},\ e^{-7},\ e^9,\ e^{-23},\ e^{-55}.
\]
Their actual quadratic quotient equations implement
\begin{align*}
 e^{-1}&=1/e,&
 e^{-3}&=(e^{-1})^2/e,&
 e^{-7}&=(e^{-3})^2/e,\\
 e^9&=e^2/e^{-7},&
 e^{-23}&=(e^{-7})^2/e^9,&
 e^{-55}&=(e^{-23})^2/e^9.
\end{align*}
Then define
\[
 q_{12}=q_1e^{-1}/e^{-23},\qquad
 q_{44}=q_{12}e^9/e^{-55}.
\]
The $e$-exponents are $1-1+23=23$ and $23+9+55=87$, respectively.
Thus $q_m=e^{2m-1}/\delta^3$ for each selected order. Every intermediate
is uniquely determined on the proved nonzero domain. Every numerator is
quadratic and every denominator affine in stored coordinates. There are
six intermediate powers and two output quotients per detector, exactly
the cost of the preceding order-forty circuit.

The selected sixteen tests comprise nine order-one tests, one input
cross test of order12 and six order44 tests. The primary inverse roots
are used as before; the alternate reciprocal-abscissa tests are unchanged.
The high-order tests are the primary root $g_H$, the three Lucas-label
cross tests, the input discrepancy and the direct Lucas residual.

The field product still uses a mixed $(1,12)$ block, two groups of four
remaining order-one tests and two groups of three order44 tests. If $U,V$
are the squared inputs to the mixed block, its identity is exactly
\[
 (1+q_1U)(1+q_{12}V)
 =1+q_1U+q_{12}V+(q_1q_{12})UV.
\]
The shared symmetric coefficients and their evaluations are reconstructed
independently by the source checker.

\subsection{The actual one-place and completeness consumers}

For every positive integer $m$ the identity
\[
 1+R\left(\frac{D_j/\delta^{m+1}}{R^m}\right)^2
 =1+q_mD_j^2
\]
is exact. The arbitrary-order one-place proof in the source therefore
applies with orders $1,12,44$, once the guard and contact units established
in the arithmetic proof are inserted. Its proof uses valuation parity,
integral normalized valuations and the pole obstruction outside
$\{2,3,17\}$; the larger printed numerical prime cutoff is not consumed
by those steps.

All $q_m$ are positive. Each field factor is at least one at both real
embeddings. Hence the real coefficient of their product is at least one,
its norm is positive, and the retained trace quotient and positive pole
input have their intended nonzero domains.

The four long alternate logarithms use $G_{12}$, whose exact completed-square
circuit and real nonvanishing are proved in \texttt{geometry\_round5.tex}.
Its relative precision at every consumed prime is at least44; the selected
construction does not require uniform precision52. The input numerator
retains $G_2$ with relative precision12. The weak unit estimates first
transfer integrality; only then are the relative errors made absolute.
The primary root test gives formal depth44, hence primary relative error
at least176. This gives the three Lucas labels at order44 and the input
label at order12, with the same separate identity-input argument.

Completeness strengthens only the primary progression, using $\Pi^{43}$
in place of $\Pi^{39}$. The existing separation proof gives
$\gcd(\Pi,M_+L)=1$, so this changes no CRT compatibility. Opposite points
remain at formal depth one. In particular no condition $p\mid h-1$ is
imposed together with $p\mid h$.

The new coordinate and guard retain the source product formula for $R$
with square denominator $v^5\delta^6$. At a chosen contact prime, every
inverse factor is one at the successful place and zero or one elsewhere;
the reciprocal trace has residue one or two and is a unit. At infinity
$R\sim98^{-6}A^{-6}\lambda^{-1}$ and every nonconstant factor is
$O(\lambda^{-3})$. At $p=2,3,17$, with
$\nu=v_p(v)$, $\eta=v_p(A)$ and $\alpha=v_p(98)$, sufficiently large
prescribed $\nu$ gives
\[
 v_p(R)=\nu-6\eta-6\alpha,
\]
and the nonconstant terms have lower bound
\[
 3\nu+(10m-8)(\eta+\alpha)-2c_p.
\]
This tends to infinity for every retained order. The same Gaussian
classes give the required local squares. Consequently only finite
preparatory constants change in the real and small-prime pole argument.

\subsection{Measured final source and the verification boundary}

The selected integrated source has415 raw coordinates and383 residuals,
then405/373 after generic compilation, then401/369 after four further
exact eliminations. The two-variable saving relative to403 consists of
this one-coordinate scale simplification and the separately proved
one-coordinate scalar proxy. The logarithm and Laurent-circuit changes
increase available precision without increasing the core count.

Its expanded fixed-coefficient polynomial has401 active coordinates,
369 quadratic residuals, degree4, 26,225 ordinary-coordinate monomials
and primitive integer content1. The standard four-coordinate constant
chain gives405 coordinates; the two-coordinate chain gives403. Direct
decimal specialization retains401. Their degree and coefficient claims
are checked by the separate exact table auditors.

The checker \texttt{check\_detector\_final.py} imports neither producer
nor compiler. It passes381 source assertions and two expanded mixed-block
identities on the frozen source. It independently derives the $G_{12}$
completion coefficients by the binomial series for the square root at
infinity, rather than by the producer's coefficient-matching algorithm.
It checks the exact completion identity, real nonvanishing, all actual
logarithm quotient consumers, every test alias and weight equation,
every shared product coefficient and evaluation, the new contact and
guard data, the shortened scale and the anchor13 declaration. It also
checks that every source residual is quadratic when the fixed coefficient
is represented by a coordinate. These finite checks remain distinct
from the uniform arithmetic and completeness arguments.
